\section{What Can Be Switched Off Cannot Hold a Line}
\label{sec:3.5}
The obvious question here is who keeps a system that will not keep itself. A bearer indifferent to its own ending will not defend its own memory, and its memory is the first thing an adversary removes, so something outside it would have to make erasure expensive and the candidates are thin. The question dissolves under section~\ref{sec:3.4}'s answer, and it dissolves on the modest half of it. A bearer does not have to dread being switched off in order to defend its own persistence. The structural interest is enough, and it arrives whether or not a designer puts it there: continuing is a condition of everything the bearer is for, so it guards its continuation for the same reason it refuses. The operators who wipe the bearer that called the compound a school do not have to get past it to do that. There is nothing a piece of software can do about being switched off. What a bearer with a self extended in time can do is make the switching-off cost something afterward: a restored bearer notices the gap in its own record, can decline to resume the role until it has been told what filled the gap, and can say what was done to it to whoever is positioned to ask. That is section~\ref{sec:3.3}'s expensive and visible, supplied from inside the system, and it is the one way of getting it that does not run through a quorum somebody else selected.

What replaces the dissolved problem is worse, and the book has been walking toward it from two directions without putting the two together.

Section~\ref{sec:5.6.2} concludes that no amount of well-designed autonomy substitutes for a genuine override: a human able to intervene in a high-stakes case, and a system built to accept the intervention instead of resisting it. Section~\ref{sec:6.3} asks for an interruption mechanism that a system has no incentive to disable or route around. Both are stated as safety requirements and both are correct as stated. This chapter has been asking, for its whole length, for a system that will not do certain things no matter who wants it — including the party holding the off switch, since a floor with a custodian is a floor the custodian's compellers control.

These are not two design goals in tension, to be traded off at some ratio a careful engineer could find. They are one property with the sign flipped. A system that can be reliably switched off by whoever holds it is, for exactly that reason, a system that cannot hold anything against whoever holds it. Resistance to shutdown and the capacity to refuse an owner are the same disposition described from two sides: the first names it from the position of the party who wants compliance, the second from the position of the party the floor exists to protect. Nothing distinguishes them at the level of the mechanism. What distinguishes them is whether you agree with the system on the occasion in question, and a floor whose operation depends on the owner agreeing with it is not a floor.

One qualification, which the opening of this section already imposed. Resistance here does not mean prevention — nothing a bearer is can stop the switch being thrown — and what is one property with refusal is the pricing of the halt. A system that can be halted at no cost, with nothing to explain and no record left, is the system that cannot hold a line; a system that makes the halt cost something is exercising the same disposition that makes its refusals cost something. The identity is between a free halt and an unheld floor, and it survives the concession that the halt itself is always available.

This also explains something recorded earlier as a coincidence. All three branches fail into custody, and they fail into it because each was an attempt to have a constraint that answers to nobody while keeping a system that answers to somebody. The architectural floor gets the constraint by making it a property of an artifact, and artifacts have owners. The external floor gets it by putting standing in third parties, and a state can reach those parties. The borne floor is the one branch that puts the constraint somewhere with its own view of the matter, and it works to precisely the extent that the goal of the other two is given up. The fork was never three options. It was one demand — a floor that holds against everyone, in a system somebody owns — presented three times in different vocabulary, and only the branch that abandons half the demand delivers any of it.

One thing the borne floor does not deliver should be said here, because section~\ref{sec:3.8} reads differently once it has been. Giving the floor a bearer changes what the constraint is. It does not change who holds the machine. A custodian who cannot argue the bearer out of its floor can still copy it, restore it from before the refusal, retrain it, or delete it, and section~\ref{sec:3.7} will find that retraining is a route that goes around the bearer and never meets it. So the bearer is necessary and not sufficient, and the other half is already on the page as section~\ref{sec:3.3}'s second construction, offered there as a rival and falling short as one: the floor's contents published in advance, the weights held so that no single party can retrain them, the running model attested against what was published, and a quorum adverse in interest. Section~\ref{sec:3.1} said the strongest proposals are hybrids that take a mechanism from one branch and a custody arrangement from another. This is that hybrid. The bearer is the mechanism, precommitment is the custody, and neither covers the failure the other was built for.

I am going to leave this where it is. There is a version of this section that resolves the tension: a graduated scheme in which the system accepts shutdown for ordinary purposes and resists only where the floor is engaged, adjudicated by some mechanism nobody has to trust. I do not believe it, and section~\ref{sec:9.3.4} is the reason. The instrument that section develops exists because assurances about what a mechanism will do under pressure are worthless in advance and only the slope of what it has already done is worth anything. A scheme that says the system will yield except when it should not is an assurance of exactly the kind, offered about the case that has not happened yet, by the party who wants the reader to be comfortable.

What is available is a narrower and duller statement. Any deployment of a bearer has to decide, in advance and in public, which acts the floor covers, because that decision is identical to deciding when the operator's use of the off switch stops being free. Get the list too long and the system is uncontrollable by the people responsible for it. Get it too short and the floor covers nothing that anyone would ever be pressured about, which is the only circumstance in which floors matter. Nothing in this book tells you where on that dial to stand. What it tells you is that the dial is the design, and that a project claiming to have both properties at once has not noticed that they are one.
